\documentclass[11pt]{article}

\usepackage{amsmath, amssymb, amsthm}
\usepackage{geometry}
\usepackage{hyperref}
\usepackage{bm}
\usepackage{booktabs}
\usepackage{xcolor}
\usepackage{listings}
\usepackage{parskip}

\geometry{margin=1.2in}

\lstset{
  basicstyle=\ttfamily\small,
  breaklines=true,
  frame=single,
  backgroundcolor=\color{gray!8},
}

\title{\textbf{Mathematical Reference}\\[0.5em]
  \large Basis Function Expansion for Gravitational Potentials\\[0.3em]
  \normalsize \texttt{bfeax} — JAX implementation}
\date{}

\begin{document}
\maketitle
\tableofcontents
\newpage

% ===========================================================================
\section{Overview}

The code converts an arbitrary density function $\rho(\mathbf{x})$ into a
smooth gravitational potential $\Phi(\mathbf{x})$ by decomposing in spherical
harmonics angularly and using the radial Green's function for Poisson's
equation.  The pipeline has four stages:

\begin{enumerate}
  \item Build a log-spaced radial grid $\{r_i\}$.
  \item At each shell $r_i$, project $\rho$ onto spherical harmonics to obtain
        the radial coefficients $\rho_{\ell m}(r_i)$.
  \item Solve the radial Poisson equation for each $(\ell,m)$ mode to obtain
        $\Phi_{\ell m}(r_i)$ using a Green's function with analytic boundary
        corrections.
  \item Fit natural cubic splines to $\Phi_{\ell m}(r)$ and $\rho_{\ell m}(r)$
        for continuous evaluation.
\end{enumerate}

The potential and reconstructed density are then evaluated as
\begin{align}
  \Phi(\mathbf{x}) &= \sum_{\ell=0}^{\ell_{\max}} \sum_{m=-\ell}^{\ell}
    \Phi_{\ell m}(r)\, Y_{\ell m}(\theta,\varphi), \\
  \rho_{\rm rec}(\mathbf{x}) &= \sum_{\ell=0}^{\ell_{\max}} \sum_{m=-\ell}^{\ell}
    \rho_{\ell m}(r)\, Y_{\ell m}(\theta,\varphi),
\end{align}
where $(r,\theta,\varphi)$ are standard spherical coordinates.  The
gravitational acceleration is $\mathbf{a} = -\nabla\Phi$, computed via
automatic differentiation.

% ===========================================================================
\section{Radial Grid}

\subsection{Log-uniform spacing}

All radial integrals are discretised on a grid of $N_r$ points uniformly
spaced in $\log r$:
\begin{equation}
  r_i = \exp\!\left(\log r_{\min} + \frac{i}{N_r - 1}
        \bigl(\log r_{\max} - \log r_{\min}\bigr)\right),
  \quad i = 0, \ldots, N_r - 1.
\end{equation}
Log spacing allocates equal resolution per decade, which is appropriate for
power-law profiles (e.g.\ NFW, Hernquist) that vary by many orders of
magnitude across the radial range.

\subsection{Parameter guidance}

\begin{center}
\begin{tabular}{lll}
\toprule
Parameter & Typical value & Notes \\
\midrule
$r_{\min}$ & $\sim 10^{-2}\,r_s$ & Well inside the inner cusp \\
$r_{\max}$ & $\geq 30\,r_s$      & Critical for NFW; outer integral diverges logarithmically \\
$N_r$      & 64--128             & 128 gives $\lesssim 0.1\%$ for NFW \\
\bottomrule
\end{tabular}
\end{center}

% ===========================================================================
\section{Real Spherical Harmonics}

\subsection{Convention}

The code uses \emph{real orthonormal} spherical harmonics.  Starting from the
complex harmonics $Y_\ell^m$, the real combinations are
\begin{equation}
  Y_{\ell m}(\theta,\varphi) =
  \begin{cases}
    \sqrt{2}\, N_{\ell m}\, P_\ell^{|m|}(\cos\theta)\,\cos(m\varphi)
      & m > 0, \\
    N_{\ell 0}\, P_\ell^0(\cos\theta)
      & m = 0, \\
    \sqrt{2}\, N_{\ell |m|}\, P_\ell^{|m|}(\cos\theta)\,\sin(|m|\varphi)
      & m < 0,
  \end{cases}
\end{equation}
with normalisation
\begin{equation}
  N_{\ell m} = \sqrt{\frac{2\ell+1}{4\pi}\frac{(\ell-|m|)!}{(\ell+|m|)!}}.
\end{equation}
These satisfy the orthonormality relation
\begin{equation}
  \int Y_{\ell m}\, Y_{\ell' m'}\,\mathrm{d}\Omega
  = \delta_{\ell\ell'}\,\delta_{mm'}.
\end{equation}

\subsection{Associated Legendre recurrence}

The associated Legendre polynomials $P_\ell^m(x)$ are computed via the
three-term recurrence (evaluated at $x = \cos\theta$).  The seed and diagonal
step are
\begin{align}
  P_0^0 &= 1, \\
  P_\ell^\ell &= -(2\ell-1)\sin\theta\, P_{\ell-1}^{\ell-1}, \\
  P_\ell^{\ell-1} &= (2\ell-1)\cos\theta\, P_{\ell-1}^{\ell-1},
\end{align}
and the upward recurrence for $m \leq \ell - 2$ is
\begin{equation}
  P_\ell^m = \frac{(2\ell-1)\cos\theta\, P_{\ell-1}^m
             - (\ell+m-1)\, P_{\ell-2}^m}{\ell - m}.
\end{equation}
All $(\ell,m)$ pairs with $0 \leq m \leq \ell \leq \ell_{\max}$ are computed
in a single sweep; the loop structure is unrolled at JAX trace time since
$\ell_{\max}$ is a static Python integer.

% ===========================================================================
\section{Angular Projection}

\subsection{Spherical harmonic decomposition at each shell}

For each radius $r_i$, the density is integrated over the unit sphere:
\begin{equation}
  \rho_{\ell m}(r_i) = \int \rho(r_i,\theta,\varphi)\, Y_{\ell m}(\theta,\varphi)\,
  \mathrm{d}\Omega.
\end{equation}
This is the fundamental step that converts the 3D density into a set of 1D
radial functions, one per mode $(\ell,m)$.

\subsection{Gauss--Legendre $\times$ uniform-$\varphi$ quadrature}

The integral over the sphere is evaluated numerically as a product quadrature.
In the $\cos\theta$ direction, $N_\theta$ Gauss--Legendre nodes $\{x_j,
w_j\}$ are used (where $x_j = \cos\theta_j$ and $w_j$ are the GL weights in
$\mathrm{d}(\cos\theta)$).  In the $\varphi$ direction, $N_\varphi$ equally
spaced nodes with weight $\Delta\varphi = 2\pi/N_\varphi$ are used.  The
discrete approximation is
\begin{equation}
  \rho_{\ell m}(r_i) \approx \sum_{j=1}^{N_\theta} \sum_{k=1}^{N_\varphi}
    \rho(r_i, \theta_j, \varphi_k)\, Y_{\ell m}(\theta_j, \varphi_k)\,
    w_j\, \Delta\varphi.
\end{equation}
Note that the GL rule integrates in $\mathrm{d}(\cos\theta)$, so the
$\sin\theta$ Jacobian factor is \emph{not} included explicitly; it is
absorbed into the GL weights.

The angular quadrature is pre-evaluated once as a tensor
$Y_{\ell m}^{jk} \equiv Y_{\ell m}(\theta_j,\varphi_k)$ of shape
$(N_{\rm modes}, N_\theta, N_\varphi)$, and the projection over all radii is
computed in a single \texttt{jnp.einsum}:
\begin{equation}
  \rho_{\ell m}(r_i) \approx \sum_{j,k} \rho(r_i,\theta_j,\varphi_k)\,
  Y_{\ell m}^{jk}\, w_j\, \Delta\varphi.
\end{equation}

\subsection{Default quadrature orders}

The Gauss--Legendre rule with $N_\theta$ points integrates polynomials of
degree $2N_\theta - 1$ exactly.  For non-polynomial integrands (e.g.\ NFW),
a factor of 3 margin is used:
\begin{equation}
  N_\theta = 3(\ell_{\max} + 2), \qquad
  N_\varphi = 4\ell_{\max} + 7 \quad (\text{kept odd to avoid aliasing}).
\end{equation}

% ===========================================================================
\section{Radial Poisson Solver}

\subsection{Green's function solution}

For each $(\ell,m)$ mode, the radial Poisson equation
\begin{equation}
  \frac{1}{r^2}\frac{\mathrm{d}}{\mathrm{d}r}\!\left(r^2
  \frac{\mathrm{d}\Phi_{\ell m}}{\mathrm{d}r}\right)
  - \frac{\ell(\ell+1)}{r^2}\Phi_{\ell m} = 4\pi G\,\rho_{\ell m}
\end{equation}
is solved via the standard Green's function (with $G=1$):
\begin{equation}
  \boxed{
  \Phi_{\ell m}(r) = -\frac{4\pi}{2\ell+1}
  \left[
    r^{-(\ell+1)} I_{\rm in}(r)
    + r^{\ell} I_{\rm out}(r)
  \right],
  }
\end{equation}
where the two cumulative integrals are
\begin{align}
  I_{\rm in}(r)  &= \int_0^r \rho_{\ell m}(r')\, r^{\prime\,\ell+2}\,
                    \mathrm{d}r', \label{eq:Iin} \\
  I_{\rm out}(r) &= \int_r^\infty \rho_{\ell m}(r')\, r^{\prime\,1-\ell}\,
                    \mathrm{d}r'. \label{eq:Iout}
\end{align}
$I_{\rm in}$ is a running integral accumulated from small $r$ outward
(``inside-out''); $I_{\rm out}$ is accumulated from large $r$ inward
(``outside-in'').  Both are evaluated by the trapezoidal rule on the
log-spaced grid.

\subsection{Trapezoidal discretisation}

Let $f_i^{\rm in} = \rho_{\ell m}(r_i)\, r_i^{\ell+2}$ and
$\Delta r_i = r_{i+1} - r_i$.  Then
\begin{equation}
  I_{\rm in}(r_i) = \Delta I_{\rm in}
  + \sum_{j=0}^{i-1} \tfrac{1}{2}\bigl(f_j^{\rm in} + f_{j+1}^{\rm in}\bigr)
    \Delta r_j,
\end{equation}
where $\Delta I_{\rm in}$ is the analytic inner tail correction (Section
\ref{sec:boundary}).  The outside-in sum for $I_{\rm out}$ is constructed
analogously via a reversed cumulative sum.

% ===========================================================================
\section{Boundary Corrections}
\label{sec:boundary}

The grid covers only $[r_{\min}, r_{\max}]$, but the integrals in
equations~\eqref{eq:Iin}--\eqref{eq:Iout} extend to $0$ and $\infty$
respectively.  Both truncations can introduce significant bias, particularly
for the NFW outer integral which diverges logarithmically unless $r_{\max}$ is
very large.

\subsection{Inner tail correction ($0 \to r_{\min}$)}

The power-law slope of $\rho_{\ell m}$ near $r_{\min}$ is estimated from the
first three grid points via a log--log finite difference:
\begin{equation}
  \alpha_{\rm in} \approx \left\langle
    \frac{\Delta\log|\rho_{\ell m}|}{\Delta\log r}
  \right\rangle_{r_0,r_1,r_2}.
\end{equation}
Assuming $\rho_{\ell m}(r) \approx A_{\rm in}\, r^{\alpha_{\rm in}}$ for
$r < r_{\min}$, the inner integral contribution is
\begin{equation}
  \Delta I_{\rm in}
  = \int_0^{r_{\min}} A_{\rm in}\, r^{\alpha_{\rm in}}\, r^{\ell+2}\,\mathrm{d}r
  = \frac{A_{\rm in}\, r_{\min}^{\alpha_{\rm in}+\ell+3}}
         {\alpha_{\rm in} + \ell + 3},
\end{equation}
which converges provided $\alpha_{\rm in} + \ell + 3 > 0$ (always satisfied
for physical density cusps with $\alpha_{\rm in} > -3$).  The correction is
added as a constant to $I_{\rm in}(r)$ for all $r \geq r_{\min}$.

The amplitude is recovered as
$A_{\rm in} = \mathrm{sign}(\rho_{\ell m}(r_0)) \cdot
|\rho_{\ell m}(r_0)| \cdot r_0^{-\alpha_{\rm in}}$,
using \texttt{exp(alpha * log(r))} to avoid overflow when $\alpha_{\rm in}$
is a traced JAX scalar.

\subsection{Outer tail correction ($r_{\max} \to \infty$)}
\label{sec:outer}

Similarly, the slope at the outer boundary is estimated from the last three
grid points:
\begin{equation}
  \alpha_{\rm out} \approx \left\langle
    \frac{\Delta\log|\rho_{\ell m}|}{\Delta\log r}
  \right\rangle_{r_{N_r-3},r_{N_r-2},r_{N_r-1}}.
\end{equation}
Assuming $\rho_{\ell m}(r) \approx A_{\rm out}\, r^{\alpha_{\rm out}}$ for
$r > r_{\max}$,
\begin{equation}
  \Delta I_{\rm out}
  = \int_{r_{\max}}^\infty A_{\rm out}\, r^{\alpha_{\rm out}}\, r^{1-\ell}\,\mathrm{d}r
  = \frac{A_{\rm out}\, r_{\max}^{2-\ell}}
         {\ell - \alpha_{\rm out} - 2},
\end{equation}
which converges provided $\alpha_{\rm out} < \ell - 2$.  For the monopole
$\ell=0$ this requires $\alpha_{\rm out} < -2$; for NFW $(\alpha_{\rm
out}\approx -3)$ this is satisfied for all $\ell \geq 0$.

\textbf{Guard conditions.}  Both corrections are applied only when
\begin{itemize}
  \item the boundary value $|\rho_{\ell m}(r_{\rm boundary})|$ exceeds
        $10^{-8}$ of the per-mode maximum (inner), or $\rho_{\ell m} > 0$
        (outer), and
  \item the power-law integral converges.
\end{itemize}
All conditionals use \texttt{jnp.where} to remain JIT-safe; neither branch
propagates NaNs.

% ===========================================================================
\section{Inner Cusp Subtraction}
\label{sec:cusp}

\subsection{Motivation}

For profiles with a steep inner cusp (e.g.\ NFW: $\rho \sim r^{-1}$), the
radial coefficients $\rho_{\ell m}(r)$ diverge toward $r=0$.  Fitting a
cubic spline directly to such data is poorly conditioned near $r_{\min}$:
small rounding errors in $\rho_{\ell m}(r_{\min})$ lead to large oscillations
in the fitted spline.

\subsection{Procedure}

For each mode $(\ell,m)$:
\begin{enumerate}
  \item Estimate the inner power-law parameters $(\alpha_{\ell m}, A_{\ell m})$
        from the first three grid points (same finite-difference estimator as
        in Section~\ref{sec:boundary}).
  \item Define the background $\mathrm{bg}(r) = A_{\ell m}\, r^{\alpha_{\ell m}}$.
  \item Fit the cubic spline to the \emph{residual}
        $\widetilde{\rho}_{\ell m}(r) = \rho_{\ell m}(r) - \mathrm{bg}(r)$,
        which is smooth and nearly flat near $r_{\min}$.
  \item At evaluation time, reconstruct
        $\rho_{\ell m}(r) = \widetilde{\rho}_{\ell m}(r) + A_{\ell m}\,
        r^{\alpha_{\ell m}}$.
\end{enumerate}

\subsection{Guard}

The subtraction is only applied when
\begin{equation}
  |\rho_{\ell m}(r_{\min})| > 10^{-6} \cdot \rho_{\rm scale},
\end{equation}
where $\rho_{\rm scale} = \max_{\ell m}\max_i |\rho_{\ell m}(r_i)|$ is the
\emph{global} amplitude scale across all modes.  This prevents spurious
large backgrounds on symmetry-zero modes (e.g.\ odd-$\ell$ modes of an
octant-symmetric density) whose values sit at machine-epsilon level.

% ===========================================================================
\section{Spline Interpolation}

\subsection{Natural cubic splines in $\log r$}

Both $\Phi_{\ell m}(r)$ and $\rho_{\ell m}(r)$ are stored as natural cubic
splines parameterised in $u = \log r$.  On each interval
$[u_i, u_{i+1}]$, the spline takes the form
\begin{equation}
  S(u) = a_i + b_i\,(u - u_i) + c_i\,(u-u_i)^2 + d_i\,(u-u_i)^3.
\end{equation}

\subsection{Fitting}

The coefficients are determined by requiring $S$ to be $C^2$ everywhere, with
natural boundary conditions $S''(u_0) = S''(u_{N_r-1}) = 0$.  This leads to
the tridiagonal system for the second derivatives $M_i = S''(u_i)$:
\begin{equation}
  h_{i-1}\,M_{i-1} + 2(h_{i-1}+h_i)\,M_i + h_i\,M_{i+1}
  = 3\!\left(\frac{y_{i+1}-y_i}{h_i} - \frac{y_i - y_{i-1}}{h_{i-1}}\right),
\end{equation}
where $h_i = u_{i+1} - u_i$ and $y_i$ are the function values.  This is
solved via \texttt{jnp.linalg.solve} on the $(N_r-2)\times(N_r-2)$ dense
matrix (safe to JIT since array shapes are static).  The polynomial
coefficients follow as
\begin{align}
  a_i &= y_i, \\
  b_i &= \frac{y_{i+1}-y_i}{h_i} - \frac{h_i\,(2M_i + M_{i+1})}{3}, \\
  c_i &= M_i, \\
  d_i &= \frac{M_{i+1} - M_i}{3\,h_i}.
\end{align}

\subsection{Batched fitting}

When building the expansion via the fast (spheroid) path, all $N_{\rm modes}$
splines share the same $\log r$ knots.  A single matrix factorisation of the
$(N_r-2)\times(N_r-2)$ system is applied to the full right-hand side matrix
of shape $(N_r-2)\times N_{\rm modes}$ simultaneously, avoiding redundant
factorisations.

\subsection{Evaluation}

For a query $u^* = \log r^*$, the interval index is found via binary search
(\texttt{jnp.searchsorted}), then the cubic polynomial is evaluated with
Horner's method.  Out-of-range queries are clamped to the nearest boundary
interval.

% ===========================================================================
\section{Potential Evaluation}

\subsection{Coordinate conversion}

Cartesian coordinates $(x,y,z)$ are converted to spherical $(r,\theta,\varphi)$:
\begin{equation}
  r = \|\mathbf{x}\|, \quad
  \theta = \arccos\!\left(\frac{z}{r}\right), \quad
  \varphi = \mathrm{atan2}(y,x) \bmod 2\pi.
\end{equation}

\subsection{Stacked evaluation path}

For expansions built via the fast path (\texttt{from\_spheroid} or
\texttt{ExpansionGrid}), all mode splines are stored as batched coefficient
arrays of shape $(N_{\rm modes}, N_r - 1)$.  A single \texttt{searchsorted}
call finds the radial interval, and the polynomial is evaluated for all modes
simultaneously:
\begin{equation}
  \Phi_{\ell m}(r^*) = a_{\ell m,\,i}
  + b_{\ell m,\,i}\,\delta u
  + c_{\ell m,\,i}\,\delta u^2
  + d_{\ell m,\,i}\,\delta u^3,
  \quad \delta u = \log r^* - u_i.
\end{equation}
The final summation is a single dot product:
\begin{equation}
  \Phi(\mathbf{x}) = \sum_{\ell m} \Phi_{\ell m}(r^*)\, Y_{\ell m}(\theta,\varphi)
  = \mathbf{\Phi}_{\rm modes} \cdot \mathbf{Y}_{\rm modes}.
\end{equation}

\subsection{Accelerations}

Gravitational accelerations are obtained via automatic differentiation:
\begin{equation}
  \mathbf{a}(\mathbf{x}) = -\nabla_{\mathbf{x}}\,\Phi(\mathbf{x}),
\end{equation}
implemented as \texttt{jax.vmap(jax.grad(\_\_call\_\_))}.  No analytical
gradient formula is required.

% ===========================================================================
\section{Scan-based Poisson Solve}

\subsection{Motivation}

Inside the JIT-compiled build functions, the naive Python loop
\texttt{for l in range(l\_max+1)} causes XLA to unroll $\ell_{\max}+1$
separate subgraphs — one per $\ell$ value — each with statically known
exponents.  For large $\ell_{\max}$ this inflates compilation time
significantly.

\subsection{Implementation}

The loop is replaced by \texttt{jax.lax.scan} over all $N_{\rm modes} =
(\ell_{\max}+1)^2$ modes in a flat ordering.  A precomputed array
$\ell_{\rm mode}[k]$ stores the $\ell$ value for the $k$-th mode.  The scan
body processes one column $\rho_{\ell m}(\cdot)$ per step:
\begin{lstlisting}
l_per_mode = [l for l in range(l_max+1) for _ in range(2*l+1)]
_, phi_T = jax.lax.scan(solve_one, None, (rho_lm_all.T, l_per_mode))
\end{lstlisting}
Because $\ell$ is now a traced float, the power-law terms are computed as
\begin{equation}
  r^{\ell+2} \;\longrightarrow\; \exp\!\bigl((\ell+2)\log r\bigr),
\end{equation}
which is numerically equivalent and avoids Python-level specialisation.  XLA
compiles a single generic body and amortises it over all steps, yielding
substantially shorter compilation time for large $\ell_{\max}$.

% ===========================================================================
\section{Convergence}

For a triaxial NFW halo ($q_1=0.8$, $q_2=0.5$) with $N_r=128$,
$r_{\max}=300\,r_s$, the density reconstruction error decreases geometrically:

\begin{center}
\begin{tabular}{cc}
\toprule
$\ell_{\max}$ & Median relative error \\
\midrule
4  & $\sim 3\%$ \\
6  & $\sim 0.8\%$ \\
8  & $\sim 0.2\%$ \\
10 & $\sim 0.06\%$ \\
\bottomrule
\end{tabular}
\end{center}

The error drops approximately $4\times$ per $\Delta\ell = 2$, consistent
with geometric (spectral) convergence in the angular expansion.  The potential
converges faster than the density; $\ell_{\max} = 2$--$4$ is often sufficient
for potential applications.

% ===========================================================================
\section{Module Summary}

\begin{center}
\begin{tabular}{ll}
\toprule
File & Responsibility \\
\midrule
\texttt{grid.py}          & Log-spaced radial grid \\
\texttt{sph\_harm.py}     & Real $Y_{\ell m}$ via Legendre recurrence; GL $\times$ uniform quadrature grid \\
\texttt{density\_coeffs.py} & Angular projection $\rho \to \rho_{\ell m}(r_i)$ via \texttt{vmap} + \texttt{einsum} \\
\texttt{spline.py}        & Natural cubic splines in JAX (JIT-safe, batched fitting) \\
\texttt{poisson.py}       & Green's function radial Poisson solver with boundary corrections \\
\texttt{potential.py}     & \texttt{MultipoleExpansion} class; fast-path \texttt{jax.lax.scan} Poisson solve \\
\bottomrule
\end{tabular}
\end{center}

\end{document}
